\documentclass[12pt]{amsart}

\usepackage[foot]{amsaddr}

\usepackage{euscript,amsfonts,amssymb,amsmath,amscd,amsthm,enumerate,hyperref}
\usepackage{mathrsfs}
\usepackage[margin=1.1in]{geometry}
\usepackage{bbm}
\usepackage{enumitem}
\usepackage[giveninits=true,style=alphabetic,backend=biber,maxnames=9,maxalphanames=9,doi=true,isbn=true,url=false,eprint=false]{biblatex}
\renewbibmacro{in:}{}
\AtEveryBibitem{
  \clearlist{language}
}
\addbibresource{Stefan.bib}
\emergencystretch=1em

\newcommand{\lb}{\left}
\newcommand{\rb}{\right}
\newcommand{\Z}{\mathbb{Z}}
\newcommand{\ev}{\mathbb{E}}
\newcommand{\Zgeq}{\mathbb{Z}_{\geq0}}
\newcommand{\Zg}{\mathbb{Z}_{>0}}
\newcommand{\R}{\mathbb{R}}
\newcommand{\Rg}{\mathbb{R}_{>0}}
\newcommand{\C}{\mathbb{C}}
\newcommand{\ind}{\mathbbm{1}}
\newcommand{\real}{\operatorname{Re}}
\newcommand{\imag}{\operatorname{Im}}
\newcommand{\dd}{\mathrm{d}}
\newcommand{\pp}{\partial}
\newcommand{\B}{\mathcal{B}}
\newcommand{\calM}{\mathcal{M}}
\newcommand{\Lip}{\mathrm{Lip}}
\newcommand{\wt}{\widetilde}
\newcommand{\wh}{\widehat}
\newcommand{\ubar}{\underline}
\newcommand{\veps}{\varepsilon}
\newcommand{\sgn}{\mathrm{sgn}}
\newcommand{\scrL}{\mathscr{L}}
\newcommand{\calG}{\mathcal{G}}
\newcommand{\calR}{\mathcal{R}}
\newcommand{\calP}{\mathcal{P}}
\newcommand{\calU}{\mathcal{U}}
\newcommand{\bF}{\mathbb{F}}
\newcommand{\bP}{\mathbb{P}}
\newcommand{\bQ}{\mathbb{Q}}
\newcommand{\E}{\mathbb{E}}
\newcommand{\M}{\mathbf{M}}
\newcommand{\Ent}{\operatorname{Ent}}
\newcommand{\Var}{\operatorname{Var}}
\newcommand{\Cov}{\operatorname{Cov}}
\newcommand{\ddt}{\frac{\dd}{\dd t}}
\newcommand{\supp}{\mathrm{supp}}
\newcommand{\Co}{\mathrm{Co}}

\theoremstyle{plain}
\newtheorem{thm}{Theorem}[section]
\newtheorem{prop}[thm]{Proposition}
\newtheorem{coro}[thm]{Corollary}
\newtheorem{lem}[thm]{Lemma}
\newtheorem{ass}[thm]{Assumption}
\newtheorem*{conj}{Conjecture}

\theoremstyle{definition}
\newtheorem{defn}[thm]{Definition}
\newtheorem{ex}[thm]{Example}
\newtheorem{cex}[thm]{Counterexample}
\newtheorem{rmk}[thm]{Remark}

\title[Sensitivity to Initial Data for the Supercooled Stefan Problem]{Sensitivity to Initial Data for Physical and Minimal Solutions of the Supercooled Stefan Problem}

\author{Graeme Baker}
\date{\today}
\address{Department of Statistics, Columbia University, New York, NY 10027, USA.}
\email{g.baker@columbia.edu}


\begin{document}

\begin{abstract}
We address the problem of well-posedness for physical and minimal solutions to a probabilistic reformulation of the supercooled Stefan problem by investigating the sensitivity of these solutions to changes in the initial data.
We show that the solution map for physical solutions is continuous under perturbations to the initial condition $X_{0-}$ (in the weak sense of probability measures), provided that $X_{0-}$ admits a unique physical solution.
Furthermore, we show continuous dependency of the solution map for minimal solutions when the data is shifted to the right; however, continuity of this solution map is shown to fail at $X_{0-}$ for shifts to the left unless uniqueness of physical solutions holds.
As a result, we show that the question of whether the solution map for minimal solutions is continuous at $X_{0-}$ is equivalent to the question of whether the given data $X_{0-}$ admits a unique physical solution.
\end{abstract}

\maketitle

\section{Introduction}

We consider the McKean--Vlasov equation
\begin{equation}\label{eq:McK-V}
\left\{
\begin{aligned}
    X_t&=X_{0-}+B_t-\Lambda_t,\quad t\ge 0\\
    \tau &= \inf \{t\ge 0: X_t\le 0\},\\
    \Lambda_t&=\alpha\bP(\tau\le t),\quad t\ge 0,
\end{aligned}
\right.
\end{equation}
where $\alpha>0$, the initial condition $X_{0-}$ is a random variable on $\R$, and $B$ is a Brownian motion independent of $X_{0-}$. In this article, we study \emph{physical} and \emph{minimal} solutions to this problem, and the sensitivity of these solutions to perturbations in the initial data. Throughout the work, we let $D(I)$ denote the space of c\`adl\`ag functions from a given interval $I$ to $\R$. \emph{A priori}, we are concerned with solutions $\Lambda\in D([0,\infty))$ to \eqref{eq:McK-V} which are in the set
\begin{align}
M=\{\ell\in D([0,\infty)):\ell_{0-}=0,\text{ and }\ell\text{ is increasing}\}.
\end{align}

The system \eqref{eq:McK-V} arises as a probabilistic reformulation of the one-phase supercooled Stefan problem in one dimension: a free boundary problem for the heat equation which models the solidification of a supercooled liquid. In that setting, the function $t\mapsto\Lambda_t$ describes the time evolution of the boundary between liquid and solid \cite{delarue_global_2022}. For general initial conditions, the corresponding free boundary partial differential equation (PDE) may exhibit finite-time blowup of the derivative of the boundary, and admit multiple weak solutions (in the sense of PDE). Equation \eqref{eq:McK-V} has also been used as a model for systemic risk in interconnected banking systems, where $\tau$ denotes the time of default for a representative bank whose assets at time $t$ are given by $X_t$, and $\Lambda_t$ is the proportion of banks which have defaulted by time $t$ (see \cite{nadtochiy_particle_2019}, as well as extensions such as \cite{nadtochiy_mean_2019,hambly_spde_2019,baker_singular_2023}, among others).
In the case of either application, equation \eqref{eq:McK-V} can be seen as the mean-field limit of the interacting particle system
\begin{equation}\label{eq:finite}
\left\{
\begin{aligned}
    X_t^{i}&=X_{0-}^i+B_t^i-L_t^N,\quad t\ge 0\\
    \tau^i &= \inf \{t\ge 0: X_t^i\le 0\},\\
    L_t^N&=\frac{\alpha }{N}\sum_{i=1}^N \ind_{\{\tau^i\le t\}} ,\quad t\ge 0,
\end{aligned}
\right.
\end{equation}
where particles $(X^i)_{1\le i\le N}$ have independent and identically distributed initial conditions $X_{0-}^i\stackrel{d}{=}X_{0-}$ and $(B^i)_{1\le i\le N}$ is a family of independent Brownian motions. 

\medskip

Depending on the initial distribution $X_{0-}$, both \eqref{eq:McK-V} and \eqref{eq:finite} may admit multiple solutions (see, for instance, \cite[Subsections 2.1 and 3.1]{delarue_particle_2015}). There is a natural way of resolving jump cascades which occur in \eqref{eq:finite} when the hitting time of one particle triggers one or more other particles to hit (for specifics, we refer the reader to the detailed discussion in \cite[Section 3]{delarue_particle_2015}). The condition used to resolve the cascades in the finite problem \eqref{eq:finite} informs which jump sizes are admissible in the limiting problem and gives rise to the following definition.

\begin{defn}
A solution $\Lambda\in M$ of \eqref{eq:McK-V} is deemed \emph{physical} provided that
\begin{align}\label{eq:phys}
\Lambda_{t}-\Lambda_{t-}=\inf\{x>0:\alpha\bP(\tau\ge t, X_{t-}\in(0,x])<x\}
\end{align}
is satisfied at all times $t\ge 0$.
\end{defn}

The physical jump condition was introduced in \cite{delarue_particle_2015} in the context of mean field networks of integrate-and-fire neurons, and later for the supercooled Stefan problem in \cite{delarue_global_2022}.
For the purpose of applications, the discontinuities of $\Lambda$ represent finite-time blowup for the free boundary in the supercooled Stefan problem, or an event where a macroscopic proportion of banks simultaneously default in the context of systemic risk. The physical jump condition \eqref{eq:phys} ensures that when a discontinuity of $\Lambda$ occurs, it is the smallest possible choice of jump among c\`adl\`ag solutions \cite[Proposition 1.2]{hambly_mckean--vlasov_2018}.

\medskip

Provided that some assumptions are made on the regularity of the initial condition $X_{0-}$, equation \eqref{eq:McK-V} is known to admit a unique physical solution. For instance, when $X_{0-}$ possesses a bounded density on $[0,\infty)$ which changes monotonicity finitely often, then the physical solution is unique and the regularity of it is fully characterized by \cite[Theorems 1.1 and 1.4]{delarue_global_2022}. Uniqueness of the physical solutions has also been established using a weaker assumption in \cite{mustapha_well-posedness_2023}, which allows for oscillatory initial densities such as those which behave like $(1+\sin 1/x)/2$ near the origin. As of writing, we do not know of an initial condition which gives rise to multiple physical solutions. Of course, we do know that such an initial condition would have to be irregular enough to invalidate the assumptions of \cite{delarue_global_2022,mustapha_well-posedness_2023}.

\medskip

With regards to the supercooled Stefan problem as a physical model, the physical solutions of \cite{delarue_global_2022} give a notion of global-in-time solutions in the presence of possible blow-ups, and these solutions are unique under the aforementioned technical assumptions.
A mathematical model for a physical phenomenon is well-posed provided that (i) there exists (appropriately defined) solutions, (ii) these solutions are unique, and (iii) they depend continuously on initial data \cite[page 7]{evans_partial_2010}. %Evans?
In Section \ref{sec:phys}, we show that the following holds
\begin{thm}
\label{thm:phys}
Let $\Lambda\in M$ solve \eqref{eq:McK-V} with initial data $X_{0-}$ such that $\E|X_{0-}|<\infty$. Let $(\Lambda^n)_{n\ge 1}$ be a sequence of physical solutions to \eqref{eq:McK-V}, where the corresponding initial conditions $(X_{0-}^n)_{n\ge 1}$ all have finite expectation. Suppose that $X_{0-}^n\to X_{0-}$ weakly and $\Lambda^n\to \Lambda$ in $(D([-1,\infty)),\operatorname{M1})$. Then $\Lambda$ satisfies
the physical jump condition
\begin{align}
\Lambda_{t}-\Lambda_{t-}=\inf\{x>0:\alpha\bP(\tau\ge t, X_{t-}\in(0,x])<x\}
\end{align}
at all $t\ge 0$. 
\end{thm}

\begin{rmk}\label{rmk:cheaptrick}
The convergence is stated in $(D([-1,\infty)),\operatorname{M1})$, where $\operatorname{M1}$ denotes the Skorokhod $\operatorname{M1}$ topology. We work on $(D([-1,\infty)),\operatorname{M1})$ rather than $(D([0,\infty)),\operatorname{M1})$ since the $\operatorname{M1}$ topology enforces continuity at the left endpoint. Throughout this work, we will use the embedding $\iota :D([0,\infty))\hookrightarrow D([-1,\infty))$ defined by $\iota f(x)=f(0-)$ for $x\in [-1,0)$ whenever $f\in D([0,\infty))$ (we will omit the symbol $\iota$). This gives us convergence at the left endpoint for free (c.f.~\cite[Remark 4.3]{cuchiero_propagation_2020}). 
\end{rmk}

With Theorem \ref{thm:phys}, we see that continuous dependence holds for the physical solution map when the initial condition $X_{0-}$ admits a unique physical solution. Therefore, we see that the question of well-posedness of physical solutions to \eqref{eq:McK-V} reduces to the question \emph{for which initial conditions does \eqref{eq:McK-V} admit a unique physical solution?} In this work, we highlight the centrality of this question for both the well-posedness of physical solutions and the well-posedness of minimal solutions, which we define next.

\medskip

Another notion of solution to equation \eqref{eq:McK-V}, known as the minimal solution and introduced in \cite{cuchiero_propagation_2020}, is arrived at as follows: first, let $\ell\in M$ be a candidate for the free boundary. Consider the problem
\begin{equation}\label{eq:Gammadef}
\left\{
\begin{aligned}
X^{\ell}_t=&X_{0-}+B_t- \ell_t,\quad t\ge 0\\
\tau^{\ell}=&\inf \{t\ge 0:X^{\ell}_t\le 0\}\\
\Gamma[\ell;X_{0-}]_t:=&\alpha\bP(\tau^{\ell}\le t), \quad t\ge 0.
\end{aligned}
\right.
\end{equation}
We see that the hunt for solutions to \eqref{eq:McK-V} amounts to finding fixed points of $\Gamma[~\cdot~;X_{0-}]$.
The operator $\Gamma[~\cdot~;X_{0-}]$ is continuous on $M$ by \cite[Proposition 2.1]{cuchiero_propagation_2020}, and it is easy to see that it is monotone in the sense that 
\begin{align}
\ell^1_t\le\ell^2_t\text{ for all }t\ge 0\implies \Gamma[\ell^1;X_{0-}]_t\le \Gamma[\ell^2;X_{0-}]_t\text{ for all }t\ge 0.
\end{align}

\medskip

For any solution $\Lambda$ of \eqref{eq:McK-V}, we must have $0\le \Lambda_t$ for all $t\ge 0$, as $\frac{1}{\alpha}\Lambda_t$ is the probability of some event. Hence
\begin{align}
\Gamma[0;X_{0-}]_t\le \Gamma[\Lambda;X_{0-}]_t=\Lambda_t,\quad t\ge 0,
\end{align}
since $\Lambda$ is a fixed point of $\Gamma[~\cdot~;X_{0-}]$. Define $\Gamma^n[~\cdot~;X_{0-}]$ recursively by
\begin{align}
\Gamma^n[~\cdot~;X_{0-}]:=\Gamma^{n-1}[\Gamma[~\cdot~;X_{0-}];X_{0-}],
\end{align}
for $n\ge 2$. Applying $\Gamma^n$ to the inequality $0\le \Lambda_t$ for all $t\ge 0$ yields 
\begin{align}\label{eq:gammaineq}
\Gamma^n[0;X_{0-}]_t\le \Gamma^n[\Lambda;X_{0-}]_t=\Lambda_t,\quad t\ge 0.
\end{align}
In \cite[Proposition 2.3]{cuchiero_propagation_2020}, it is shown that the limit $\ubar{\Lambda}:=\lim_{n\to \infty}\Gamma^n[0;X_{0-}]$ is a solution to \eqref{eq:McK-V}.
\begin{defn}
We call $\ubar{\Lambda}:=\lim_{n\to \infty}\Gamma^n[0;X_{0-}]$ the \emph{minimal} solution to \eqref{eq:McK-V}.
\end{defn}
Taking the limit of \eqref{eq:gammaineq}, we see that $\ubar{\Lambda}_t\le\Lambda_t$, for all $t\ge 0$ when $\Lambda$ is any other solution to \eqref{eq:McK-V}.
The minimality of $\ubar{\Lambda}$ among solutions to \eqref{eq:McK-V} can also be taken as the definition of the minimal solution. Either by construction, or by minimality, we have that the minimal solution is unique. We note the following important connection between minimal and physical solutions:
\begin{prop}[Theorem 6.5 in \cite{cuchiero_propagation_2020}]\label{prop:minisphys}
When $\E|X_{0-}|<\infty$, the minimal solution is a physical solution.
\end{prop}
From \cite{cuchiero_propagation_2020}, there is the following conjecture concerning sensitivity to initial conditions for minimal solutions:
\begin{conj}[Conjecture 6.10 in \cite{cuchiero_propagation_2020}]
Let $X_{0-}$ be given with $\E|X_{0-}|<\infty$. For $x\in\R$, let $\ubar{\Lambda}^x$ denote the minimal solution to \eqref{eq:McK-V} with shifted initial condition $X_{0-}+x$. The map 
$x\mapsto \ubar{\Lambda}^x$ is continuous from $\R$ to $(D([-1,\infty)),\operatorname{M1})$.
\end{conj}

In a positive direction, we are able to show that $x\mapsto \ubar{\Lambda}^x$ is c\`adl\`ag in the following sense:

\begin{thm}
Suppose that $\E|X_{0-}|<\infty$. Then $x\mapsto \ubar{\Lambda}^x$ is right-continuous from $\R$ to $(D([-1,\infty)),\operatorname{M1})$. As $x\uparrow 0$, $\ubar{\Lambda}^x$ converges in $(D([-1,\infty)),\operatorname{M1})$ to some $\Lambda^0\in M$, which is a solution to \eqref{eq:McK-V}.
\end{thm}
\begin{proof}
Right continuity is a consequence of our more general convergence result, Theorem \ref{thm:rc}. The proof of that theorem shows that the left limit gives some solution to \eqref{eq:McK-V}. This is formally stated in Corollary \ref{coro:ll}.
\end{proof}
Returning to the notion of well-posedness, the conjecture effectively states that the problem of finding minimal solutions to \eqref{eq:McK-V} is well-posed among initial conditions with $\E|X_{0-}|<\infty$. 
Proposition \ref{prop:minisphys} tells us that the minimal solutions $\ubar{\Lambda}^x$ are physical, and hence by Theorem \ref{thm:phys}, the left limit point $\Lambda^0$ is physical. Therefore, if the initial condition $X_{0-}$ admits a unique physical solution, then $\Lambda^0=\ubar{\Lambda}^0$ and the conjecture holds (this was previously known, see \cite[Theorem 6.6]{cuchiero_propagation_2020}). 
To complicate matters, the following theorem shows that the non-uniqueness of physical solutions to \eqref{eq:McK-V} implies failure of left-continuity for the map $x\mapsto\ubar{\Lambda}^x$. In fact, the left limit $\Lambda^0=\lim_{x\uparrow 0}\ubar{\Lambda}^x$ is a \emph{maximal} physical solution in that it dominates all other physical solutions at all times.

\begin{thm}
\label{thm:nonunique}
Let $X_{0-}$ be an initial condition with $\E|X_{0-}|<\infty$ and let $\Lambda^0=\lim_{x\uparrow 0}\ubar{\Lambda}^x$. If $\Lambda$ is any physical solution to \eqref{eq:McK-V} with initial condition $X_{0-}$ then $\Lambda^0_t\ge \Lambda_t$ for all $t\ge 0$. As a result, if $X_{0-}$ admits at least two physical solutions, then $\Lambda^0\neq\ubar{\Lambda}^0$.
\end{thm}

Therefore, the question of well-posedness for the minimal solutions to \eqref{eq:McK-V}, just as with well-posedness for physical solutions, reduces to the question \emph{for which initial conditions does \eqref{eq:McK-V} admit a unique physical solution?} In other words, the conjecture holds if and only if all initial conditions with finite expectation admit unique physical solutions.

\medskip

Our general convergence result for minimal solutions, stated below, allows for both deterministic and random shifts to the initial condition: we consider a sequence of initial conditions $(X_{0-}^n)_{n\ge 1}$ and assume that the cumulative distribution functions (CDFs) of the initial conditions increase pointwise and converge to the CDF for $X_{0-}$ at all continuity points of the latter. 
Note that pointwise convergence of CDFs at continuity points of the limit is equivalent to weak convergence of probability measures on $\R$, by the Helly-Bray Theorem (see, for instance, \cite[Theorem 11.1.2]{dudley_real_2002}). Alternatively, we may consider $X_{0-}^n=X_{0-}+Y^n$ for a sequence of independent non-negative random variables $(Y^n)_{n\ge 1}$ with $Y^n\to0$ weakly and $\bP(Y^n\le x)\ge \bP(Y^m\le x)$ for all $x>0$ whenever $n>m$ (we note that this partial ordering is known as stochastic dominance).

\begin{thm}
\label{thm:rc}
Let $X_{0-}$ have CDF $F$. Let $(X_{0-}^n)_{n\ge 1}$ be random variables on $[0,\infty)$ with CDFs $(F^n)_{n\ge 0}$, respectively, such that $F(x) \ge F^n(x)\ge F^m(x)$ for all $x\ge 0$ whenever $m<n$, and $F^n(x)\to F(x)$ whenever $x$ is a continuity point of $F$. For $n\ge0$, let $\ubar{\Lambda}^n$ denote the minimal solution to \eqref{eq:McK-V} with initial data $X_{0-}^n$; and let $\ubar{\Lambda}$ denote the minimal solution with initial data $X_{0-}$. Then $\ubar{\Lambda}^n\to \ubar{\Lambda}$ in $(D([-1,\infty)),\operatorname{M1})$
\end{thm}

In \cite[Proposition 5.1]{hambly_mckean-vlasov_2022}, convergence for the initial conditions $X_{0-}^n=X_{0-}+\xi^n$ is considered, where $\xi^n$ is an exponential random variable with rate $n$. Following the reasoning of \cite[Section 3]{hambly_mckean-vlasov_2022}, we see that this in turn is a generalization of our earlier result \cite[Theorem 1.2]{baker_zero_2022}, where it was assumed that the density for $X_{0-}$ is bounded by $\alpha/2$ 
%mention Misha too?
(this bound renders the physical solution unique and continuous by \cite[Theorem 2.2 and last paragraph of Section 2.1]{ledger_uniqueness_2020}).
Our proof of Theorem \ref{thm:rc} fixes an error in the proof of \cite[Proposition 5.1]{hambly_mckean-vlasov_2022} (see Remark \ref{rmk:error}) and further generalizes the result by using a fact about the zeroes of Brownian motion with a bounded variation drift. Specifically, we make use of \cite[proof of Proposition 3.4(i)]{antunovic_isolated_2011} to show that almost surely none of these zeroes are local maxima. This observation has also found recent use in the works \cite{baker_singular_2023,nadtochiy_stefan_2022}.

\medskip

The outline of the remainder of the paper is as follows: in the following subsection we will cover some technical preliminaries. Then, in Sections \ref{sec:rc} and \ref{sec:phys} we give the proofs of our convergence results for minimal and physical solutions, respectively. Finally, in Section \ref{sec:nonunique}, we prove Theorem \ref{thm:nonunique}, thus showing that non-uniqueness of physical solutions implies the ill-posedness of minimal solutions.

\subsection{Mathematical Preliminaries}

We first state a useful characterization for convergence of monotone functions in the M1 topology. 

\begin{prop}[Corollary 12.5.1 in \cite{whitt_stochastic-process_2002}]\label{prop:M1conv}
Let $D([T_0,T_1])$ denote the space of c\`adl\`ag functions from $[T_0,T_1]$ to $\R$ which are left continuous at $T_0$. Suppose that $(f^n)_{n\ge 1}$ is a family of non-decreasing functions in $D([T_0,T_1])$. Then $f^n\to f$ in $(D([T_0,T_1]),\operatorname{M1})$ if and only if $f^n(t)\to f(t)$ for all $t$ in a dense subset of $[T_0,T_1]$ including $T_0$ and $T_1$.
\end{prop}

Following \cite{whitt_stochastic-process_2002}, we say that $f^n\to f$ in $(D([T_0,\infty)),\operatorname{M1})$ if there is a sequence $(T_N)_{N\ge1}$ with $T_N\nearrow\infty$ such that $f^n\to f$ in $(D([T_0,T_1]),\operatorname{M1})$ for all $N\ge 1$. Note that this imposes convergence at the left endpoint: $f^n(T_0)\to f(T_0)$. 
Recall from Remark \ref{rmk:cheaptrick} that we are working with convergence in $(D([-1,\infty)),\operatorname{M1})$, and we will embed any function $f\in D([0,\infty))$ into $D([-1,\infty))$ by defining $f(x)=f(0-)$ for $x\in [-1,0)$.

\medskip

For the proof of Theorem \ref{thm:rc}, we will also make use of the following definition and technical lemma from \cite[Definition 5.1 and Lemma 5.4]{cuchiero_propagation_2020}.

\begin{defn}
For $f\in D([-1,\infty))$ define the path functionals
\begin{align}
    \tau_0(f):=\inf\{s\ge 0: f_s\le 0\}\text{ and }\lambda_t(f):=\ind_{\{\tau_0(f)\le t\}}.
\end{align}
\end{defn}

\begin{lem}\label{lem:cross}
Let $(f,f^1,f^2,\dots)\subset D([0,\infty))$. Suppose that $f^n\to f$ in $(D([-1,\infty)),\operatorname{M1})$ and $f$ satisfies the crossing property
\begin{align}\label{eq:cross}
\inf_{0\le s \le h}(f_{\tau_0(f)+s}-f_{\tau_0(f)})<0\text{ for all }h>0.
\end{align}
Then there exists some co-countable set $E$ such that
\begin{align}
\lim_{n\to\infty} \lambda_t(f^n)=\lambda_t(f)
\end{align}
for all $t\in E$.
\end{lem}

\section{Convergence of Minimal Solutions}\label{sec:rc}

We now present the proof of our general convergence result for minimal solutions.

\begin{proof}[Proof of Theorem \ref{thm:rc}]
The condition $F(x)\ge F^n(x)\ge F^m(x)$ for all $x\ge 0$ whenever $m<n$ implies
\begin{align}
\Gamma^{k}[0;X_{0-}^m]_t\le \Gamma^{k}[0;X_{0-}^n]_t \le \Gamma^{k}[0;X_{0-}]_t,\quad t\ge 0
\end{align}
for all $k\ge 1$. Taking the limit as $k\to\infty$ yields $\ubar{\Lambda}^m_t\le \ubar{\Lambda}^n_t \le \ubar{\Lambda}_t$ for all $t\ge 0$. For each $t\ge 0$ we define the pointwise limit
\begin{align}\label{eq:pw}
\Lambda^\infty_t := \lim_{n\to\infty} \ubar{\Lambda}^n_t
\end{align}
and note that $\Lambda^\infty_t\le  \ubar{\Lambda}_t$.

\medskip

The function $t\mapsto \Lambda^\infty_t$ is non-decreasing and hence has at most countably many discontinuities, which we denote by $J$. $\Lambda^\infty$ may not be c\`adl\`ag, so we define $\wh{\Lambda}^\infty_t = \lim_{s\downarrow t}\Lambda^\infty_s$. We have that $\ubar{\Lambda}^n_t \to \wh{\Lambda}^\infty_t$ whenever $t\in J^c$. We extend the domain of all relevant functions to $[-1,\infty)$ by setting them to zero on all of $[-1,0)$. Proposition \ref{prop:M1conv}, gives $\ubar{\Lambda}^n\to \wh{\Lambda}^\infty$ in $(D([-1,t]),\operatorname{M1})$ for all $t\in J^c$ and hence on $(D([-1,\infty)),\operatorname{M1})$.

\medskip

Provided we can show
\begin{align}\label{eq:limit}
\wh\Lambda^\infty_t=\bP (\inf_{0\le s\le t}X_{0-}+B_s-\alpha\wh\Lambda^\infty_s\le 0),
\end{align}
we will have $\wh\Lambda^\infty = \ubar{\Lambda}$ by minimality of the latter. Since $\wh\Lambda^\infty$ is increasing and bounded, \cite[proof of Proposition 3.4(i)]{antunovic_isolated_2011} implies that $t\mapsto X_{0-}+B_t-\wh\Lambda^\infty_t$ satisfies the crossing property \eqref{eq:cross}.
Recall that the pointwise convergence of $F^n\to F$ at continuity points implies that $X^n_{0-}\to X_{0-}$ in law, by the Helly-Bray Theorem \cite[Theorem 11.1.2]{dudley_real_2002}.
Therefore, we may apply Lemma \ref{lem:cross} to obtain
\begin{align}\label{eq:verified}
\lim_{n\to\infty}\bP (\inf_{0\le s\le t}X_{0-}+B_s-\alpha\ubar{\Lambda}^n_s\le 0)=\bP (\inf_{0\le s\le t}X_{0-}+B_s-\alpha\wh\Lambda^\infty_s \le 0)
\end{align}
for $t$ in a co-countable set $E$. The left hand side of the above expression is the definition of $\wh\Lambda^\infty_t$ when $t\in J^c$.
Hence $\wh\Lambda^\infty_t = \ubar{\Lambda}_t$ for $J^c \cap E$ (and in fact everywhere by the right continuity of both) and $\ubar{\Lambda}^n\to \ubar{\Lambda}$ in $(D([-1,\infty)),\operatorname{M1})$. \qedhere

\end{proof}

\begin{rmk}\label{rmk:error}
We note an error in the proof of \cite[Proposition 5.1]{hambly_mckean-vlasov_2022}: the last display equation in that proof gives
\begin{align}
\lim_{\kappa\to \infty}\Lambda^\kappa_t=\lim_{\kappa\to \infty}\bP (X_{0-}+\xi^\kappa +B_t-\alpha\Lambda^\kappa_t\le 0)=\bP (X_{0-}+B_t-\alpha\Lambda^\infty_t\le 0),
\end{align}
which is similar to \eqref{eq:verified}, except that the infimum is missing from the expression for $\Lambda^\kappa_t$. This equality is then verified by the Portmanteau Theorem (see, e.g., \cite[11.1.1]{dudley_real_2002}) since $X_{0-}+B_t-\alpha\Lambda^\infty_t$ has no atoms. However, this fails when the infimum is present. 
\end{rmk}

The proof of Theorem \ref{thm:rc} also gives us some information about left limits: consider a sequence $X^n_{0-}\to X_{0-}$ with $F\le F^n\le F^m$ whenever $m<n$. This holds, for instance, when $X^n_{0-}=X_{0-}-x_n$ where $(x_n)_{n\ge 1}$ is a sequence of non-negative numbers converging monotonically to zero. Then $\ubar{\Lambda}^m_t\ge \ubar{\Lambda}^n_t \ge \ubar{\Lambda}_t$ for all $t\ge 0$ and one may define a pointwise limit $\Lambda^\infty$ analogous to \eqref{eq:pw}. Continuing through the proof we see that the right continuous version $\wh\Lambda^\infty$ solves \eqref{eq:McK-V}. However, without the inequality $\Lambda^\infty_t\le \ubar{\Lambda}_t$ at all $t\ge 0$, we cannot conclude that $\wh\Lambda^\infty$ is a minimal solution. We summarize the considerations of this paragraph as a corollary. 
\begin{coro}\label{coro:ll}
Let $X_{0-}$ have CDF $F$. Let $(X_{0-}^n)_{n\ge 1}$ be random variables on $[0,\infty)$ with CDFs $(F^n)_{n\ge 0}$, respectively, such that $F(x) \le F^n(x)\le F^m(x)$ for all $x\ge 0$ whenever $m<n$, and $F^n(x)\to F(x)$ whenever $x$ is a continuity point of $F$. For $n\ge0$, let $\ubar{\Lambda}^n$ denote the minimal solution to \eqref{eq:McK-V} with initial data $X_{0-}^n$. Then $\Lambda^{\infty}$ defined by
\begin{align}
\Lambda^\infty_t = \lim_{s\downarrow t}\lim_{n\to\infty}\ubar{\Lambda}^n_s
\end{align}
solves \eqref{eq:McK-V} with initial data $X_{0-}$.
\end{coro}

\section{Convergence of Physical Solutions}\label{sec:phys}

In this section, we show that physical solutions converge to solutions which satisfy the physical jump condition. Our proof is by a similar method to \cite[Proof of Theorem 6.4]{cuchiero_propagation_2020}, except that in our case the limit is being approximated by solutions to \eqref{eq:McK-V} with converging initial conditions and free boundaries rather than by finite particle systems. Note also that $\alpha$ appears in a different location in that work.

\begin{proof}[Proof of Theorem \ref{thm:phys}]
Fix $T>0$. For $n \ge 1$ and $t \ge 0$, we let $X_t^n=X_{0-}^n+B_t- \Lambda^n_t$, $\tau^n=\inf \{t\ge 0:X_t^n\le 0\}$, and consider the subprobability measures $\nu_{t-}^n$ defined by 
\begin{align}
\nu_{t-}^n(A)=\bP (\tau^n \ge t, X_t^n\in A).
\end{align}
on $t\in [0,T]$, and $\nu_{t-}^n=0$ on $[-1,0)$.
We define $\nu_{t-}$ analogously for $X_t=X_{0-}+B_t- \Lambda_t$ and $\tau=\inf \{t\ge 0:X_t\le 0\}$, which along with $\Lambda$ give a solution to \eqref{eq:McK-V} by the reasoning in the proof of Theorem \ref{thm:rc}.
 Each of the bounded increasing functions $\Lambda, \Lambda_0, \Lambda_1,\dots$ can have at most countably many discontinuities, and we collect all of these together in a set $J$. Recalling Proposition \ref{prop:M1conv}, we have that $\Lambda^n_t\to\Lambda_t$ for all $t$ in a dense set, which we will call $E$.
 
\medskip

By \cite[Proof of Theorem 6.4, Step 2]{cuchiero_propagation_2020}, for sufficiently small $\veps>0$ for every $n\ge 1$ we have that
\begin{align}
\alpha\nu_{t-}^n([0,z+C\veps^{1/3}])\ge z(1-C\veps^{1/3})
\end{align}
whenever $t,t+\veps\in[0,T]\cap J^c\cap E$ and $z<\Lambda^n_{t+\veps}-\Lambda^n_{t}-C\veps^{1/3}$. The constant $C$ depends only on $\alpha$, $\E|X_{0-}|$, and $T$. 
Using the weak convergence of $X_{0-}^n$ to $X_{0-}$, the $\operatorname{M1}$ convergence of $\Lambda^n$ to $\Lambda$, and following \cite[Proof of Theorem 6.4, Step 1]{cuchiero_propagation_2020} we obtain
\begin{align}\label{eq:subprobn}
\lim_{n\to\infty}\nu_{t-}^n=\nu_{t-}
\end{align}
%in what sense
weakly at all $t\in [-1,T]\cap J^c\cap E$. Note also that 
\begin{align}\label{eq:subprobtime}
\lim_{s\uparrow t}\nu_{s-}^n=\nu_{t-}
\end{align}
by \cite[Proof of Theorem 6.4, first argument of Step 2]{cuchiero_propagation_2020}.

\medskip

Suppose now that $\Lambda_t-\Lambda_{t-}=L>0$ for some $t\in[0,T)$. Then, using the density of $J^c\cap E$ in $[-1,T]$, we may find $(t_m)_{m\ge 1}$ and $(\veps_m)_{m\ge 1}$ with $t_m,t_m+\veps_m\in J^c\cap E$ such that $t_m\uparrow t$, $t_m+\veps_m>t$, and $\veps_m\downarrow 0$. Let $z<L$. Since $\Lambda$ is c\`adl\`ag we must have $z<\Lambda_{t_m+\veps_m}-\Lambda_{t_m}-C\veps_m^{1/3}$ for sufficiently large $m$. Since $t_m,t_m+\veps_m\in J^c\cap E$, then by the definition of $E$, for any $\eta>0$ there is some $N$ such that $\Lambda_{t_m+\veps_m}-\Lambda_{t_m}-C\veps_m^{1/3}<\Lambda^n_{t+\veps_m}-\Lambda^n_{t_m}-C\veps_m^{1/3}+\eta$ for all $n\ge N$. For any $\delta>0$, taking $\eta$ smaller if necessary, using the Portmanteau theorem, along with \eqref{eq:subprobn} and \eqref{eq:subprobtime} gives
\begin{align}
\alpha\nu_{t-}([0, z+\delta])\ge \limsup_{m\to\infty}\limsup_{n\to\infty} \alpha\nu_{t_m}^n([0, z+\delta])\ge \limsup_{m\to\infty} z(1-C\veps_m^{1/3})=z.
\end{align}
Furthermore, by the argument of \cite[Proof of
Proposition 1.2]{hambly_mckean--vlasov_2018}, we must have $L\ge \inf\{x>0:\alpha\bP(\tau\ge t, X_{t-}\in(0,x])<x\}$ since $\Lambda$ is right continuous.  
\end{proof}

\section{Non-uniqueness of Physical Solutions Implies Ill-posedness of Minimal Solutions}\label{sec:nonunique}

In this final section we prove Theorem \ref{thm:nonunique}, hence showing that the minimal solution map at a given initial condition is continuous if and only if the initial condition admits a unique physical solution.

\begin{proof}[Proof of Theorem \ref{thm:nonunique}]
Let $x<0$ and let $\ubar{\Lambda}^x$ be the minimal solution to \eqref{eq:McK-V} with initial condition $X_{0-}+x$.
By Proposition \ref{prop:minisphys}, $\ubar{\Lambda}^x$ is physical. At time zero, the physical jump condition \eqref{eq:phys} yields
\begin{equation}
\begin{split}
\ubar{\Lambda}^x_0-\ubar{\Lambda}^x_{0-}&=\inf\{z>0:\alpha\bP(X_{0-}+x\in(0,z])<z\}\\
&\ge \inf\{z>0:\alpha\bP(X_{0-}\in(0,z])<z\}\\
&=\Lambda_{0}-\Lambda_{0-}.
\end{split}
\end{equation}
Consider 
\begin{align}
t^*:=\inf\{t\ge0:\ubar{\Lambda}^x_t<\Lambda_t\}
\end{align}
and suppose for the sake of contradiction that $t^*<\infty$.
First, we consider the case where 
\begin{align}\label{eq:t*iszero}
t^*=0
\end{align}
and show that this is impossible. By right-continuity of $\Lambda$, for any $\veps>0$ we can find a $\delta>0$ such that $\Lambda_t-\Lambda_0<\veps$ for all $t<\delta$. Assuming \eqref{eq:t*iszero} holds, we may always find some $t_\veps<\delta$ such that
\begin{align}\label{eq:tepsineq}
\ubar{\Lambda}^x_{t_\veps}<\Lambda_{t_\veps}.
\end{align}
On $[0,t_\veps]$, we may upper bound $\Lambda_t-\Lambda_0$ by $\veps$ and lower bound $\ubar{\Lambda}^x_t-\ubar{\Lambda}^x_0$ by zero. Taking $\veps<|x|$ and applying \eqref{eq:McK-V} gives
\begin{equation}
\begin{split}
\ubar{\Lambda}^x_{t_\veps}&\ge \alpha\bP \lb(\inf_{0\le s\le t_\veps} X_{0-}+x+B_s-(\ubar{\Lambda}^x_0-\ubar{\Lambda}^x_{0-})\le 0\rb)\\
&\ge \alpha\bP \lb(\inf_{0\le s\le t_\veps} X_{0-}+B_s-(\Lambda_0-\Lambda_{0-}+\veps)\le 0\rb)\\
&\ge \Lambda_{t_\veps},
\end{split}
\end{equation}
which contradicts \eqref{eq:tepsineq}. Therefore, we may assume that $t^*>0$.

\medskip

The goal now is to show that $\ubar{\Lambda}^x_{t^*}>\Lambda_{t^*}$, which along with the definition of $t^{*}$ will contradict the right-continuity of $\ubar{\Lambda}^x$ and $\Lambda$.
Suppose that $\ubar{\Lambda}^x_{t^*}\le\Lambda_{t^*}$ and let \newline ${\tau=\inf\{t\ge 0: X_{0-}+B_t-\Lambda_t\le 0\}}$. Then by the definition of $\Lambda$ in \eqref{eq:McK-V} we must have
\begin{equation}\label{eq:rcfail}
\begin{split}
\Lambda_{t^*-}+(\ubar{\Lambda}^x_{t^*}-\Lambda_{t^*-})&\le\alpha\bP \lb(\inf_{0\le s< t^*} X_{0-}+B_s-\Lambda_{s}\le 0\rb)\\
 &\quad\quad+\alpha\bP \bigg(\tau\ge t^*, X_{0-}+B_{t^*}\in[\Lambda_{t*-},\ubar{\Lambda}^x_{t^*}]\bigg),
\end{split}
\end{equation}
which holds with equality when $\ubar{\Lambda}^x_{t^*}=\Lambda_{t^*}$, and with possible inequality when $\ubar{\Lambda}^x_{t^*}<\Lambda_{t^*}$ by applying the physical jump condition \eqref{eq:phys} to $\Lambda$ at time $t^*-$. Consider the following augmented path defined on $[0,t^*]$:
\begin{align}
\wt{\Lambda}_t=
\begin{cases}
\Lambda_t, & t\in [0,t^*),\\
\ubar{\Lambda}^x_{t^*}, & t=t^*.
\end{cases}
\end{align}
We see that the right hand side of \eqref{eq:rcfail} is equal to 
\begin{align}
\alpha\bP \lb(\inf_{0\le s\le t^*} X_{0-}+B_s-\wt\Lambda_{s}\le 0\rb),
\end{align}
where we note that the infimum here is taken over $[0,t^*]$ rather than $[0,t^*)$.
Note that, ${\ubar{\Lambda}^x_t\ge\Lambda_t}$ on $[0,t^*)$ by the definition of $t^*$. Therefore, $\ubar{\Lambda}^x_t\ge\wt{\Lambda}_t$ on $[0,t^*]$. Since $|x|>0$ and $t^*>0$, the following holds with strict inequality:
\begin{align}
\alpha\bP \lb(\inf_{0\le s\le t^*} X_{0-}+B_s-\wt\Lambda_{s}\le 0\rb) < \alpha\bP \lb(\inf_{0\le s\le t^*} X_{0-}+x+B_s-\ubar{\Lambda}^x_{s}\le 0\rb)=\ubar{\Lambda}^x_{t^*}.
\end{align}
Putting this together with \eqref{eq:rcfail} yields $\ubar{\Lambda}^x_{t^*}<\ubar{\Lambda}^x_{t^*}$. This contradiction implies that ${\ubar{\Lambda}^x_{t^*}>\Lambda_{t^*}}$, which as mentioned before, contradicts the right-continuity of $\ubar{\Lambda}^x$ and $\Lambda$ at $t^*$. Hence ${t^*=\infty}$ and we must have $\ubar{\Lambda}^x_t\ge\Lambda_t$ for all $t\ge0$. As $x$ was arbitrary, we may take the limit as $x$ goes to 0, and the inequality $\Lambda^0_t\ge\Lambda_t$ holds for all $t\ge 0$.

\medskip

Suppose now that $X_{0-}$ admits two non-equal physical solutions $\Lambda$ and $\hat{\Lambda}$. Suppose there exists some $t>0$ such that $\Lambda_t>\hat{\Lambda}_t$ (if this is not the case then we can swap the labelling on $\Lambda$ and $\hat{\Lambda}$ so that this holds). Then by the above arguments, ${\Lambda^0_t\ge\Lambda_t>\hat{\Lambda}_t\ge \ubar{\Lambda}_t}$. This shows that $\Lambda^0\neq \ubar{\Lambda}$. \qedhere

\end{proof}

\section*{Acknowledgements}

The author was partially supported by an NSERC PGS-D scholarship and a Princeton SEAS innovation research grant. The author extends gratitude to his PhD advisor, Misha Shkolnikov; as well as Hezekiah Grayer II, Scander Mustapha, and Stefan Rigger for many useful conversations.

\printbibliography

\bigskip\bigskip

\end{document}






% $ biblatex auxiliary file $
% $ biblatex bbl format version 3.0 $
% Do not modify the above lines!
%
% This is an auxiliary file used by the 'biblatex' package.
% This file may safely be deleted. It will be recreated by
% biber as required.
%
\begingroup
\makeatletter
\@ifundefined{ver@biblatex.sty}
  {\@latex@error
     {Missing 'biblatex' package}
     {The bibliography requires the 'biblatex' package.}
      \aftergroup\endinput}
  {}
\endgroup


\refsection{0}
  \datalist[entry]{anyt/global//global/global}
    \entry{antunovic_isolated_2011}{article}{}
      \name{author}{4}{}{%
        {{hash=bded24eba68a3269334bea34d3344999}{%
           family={Antunovic},
           familyi={A\bibinitperiod},
           given={Tonci},
           giveni={T\bibinitperiod}}}%
        {{hash=5f48a88a8f20e2a5bc17bf5070aa179c}{%
           family={Burdzy},
           familyi={B\bibinitperiod},
           given={Krzysztof},
           giveni={K\bibinitperiod}}}%
        {{hash=be791fcf2bac698f7a6b715cdf5d5d51}{%
           family={Peres},
           familyi={P\bibinitperiod},
           given={Yuval},
           giveni={Y\bibinitperiod}}}%
        {{hash=d255f4f3bb84334800ae11d7cebebe1a}{%
           family={Ruscher},
           familyi={R\bibinitperiod},
           given={Julia},
           giveni={J\bibinitperiod}}}%
      }
      \list{language}{1}{%
        {en}%
      }
      \strng{namehash}{31d68386eb219c093ae92731f7d0ebb3}
      \strng{fullhash}{31d68386eb219c093ae92731f7d0ebb3}
      \strng{bibnamehash}{31d68386eb219c093ae92731f7d0ebb3}
      \strng{authorbibnamehash}{31d68386eb219c093ae92731f7d0ebb3}
      \strng{authornamehash}{31d68386eb219c093ae92731f7d0ebb3}
      \strng{authorfullhash}{31d68386eb219c093ae92731f7d0ebb3}
      \field{labelalpha}{ABPR11}
      \field{sortinit}{A}
      \field{sortinithash}{d77c7cdd82ff690d4c3ef13216f92f0b}
      \field{labelnamesource}{author}
      \field{labeltitlesource}{title}
      \field{abstract}{It is well known that standard one-dimensional Brownian motion B(t) has no isolated zeros almost surely. We show that for any α {<} 1/2 there are α-Hölder continuous functions f for which the process B − f has isolated zeros with positive probability. We also prove that for any continuous function f , the zero set of B − f has Hausdorff dimension at least 1/2 with positive probability, and 1/2 is an upper bound on the Hausdorff dimension if f is 1/2-Hölder continuous or of bounded variation .}
      \field{issn}{1083-6489}
      \field{journaltitle}{Electronic Journal of Probability}
      \field{month}{1}
      \field{number}{65}
      \field{title}{Isolated {Zeros} for {Brownian} {Motion} with {Variable} {Drift}}
      \field{urlday}{20}
      \field{urlmonth}{4}
      \field{urlyear}{2021}
      \field{volume}{16}
      \field{year}{2011}
      \field{urldateera}{ce}
      \verb{doi}
      \verb 10.1214/EJP.v16-927
      \endverb
      \verb{file}
      \verb Antunovic et al. - 2011 - Isolated Zeros for Brownian Motion with Variable D.pdf:/Users/graemeb/Zotero/storage/MRDKCSWU/Antunovic et al. - 2011 - Isolated Zeros for Brownian Motion with Variable D.pdf:application/pdf
      \endverb
      \verb{urlraw}
      \verb https://projecteuclid.org/journals/electronic-journal-of-probability/volume-16/issue-none/Isolated-Zeros-for-Brownian-Motion-with-Variable-Drift/10.1214/EJP.v16-927.full
      \endverb
      \verb{url}
      \verb https://projecteuclid.org/journals/electronic-journal-of-probability/volume-16/issue-none/Isolated-Zeros-for-Brownian-Motion-with-Variable-Drift/10.1214/EJP.v16-927.full
      \endverb
    \endentry
    \entry{baker_zero_2022}{article}{}
      \name{author}{2}{}{%
        {{hash=0761dfc2a74ec6f0d1effbf27f715d59}{%
           family={Baker},
           familyi={B\bibinitperiod},
           given={Graeme},
           giveni={G\bibinitperiod}}}%
        {{hash=cce28fd978c3edebbc2aee9c6f7e9b0f}{%
           family={Shkolnikov},
           familyi={S\bibinitperiod},
           given={Mykhaylo},
           giveni={M\bibinitperiod}}}%
      }
      \list{language}{1}{%
        {en}%
      }
      \strng{namehash}{5438078037fd2ae6b07b804b456fc748}
      \strng{fullhash}{5438078037fd2ae6b07b804b456fc748}
      \strng{bibnamehash}{5438078037fd2ae6b07b804b456fc748}
      \strng{authorbibnamehash}{5438078037fd2ae6b07b804b456fc748}
      \strng{authornamehash}{5438078037fd2ae6b07b804b456fc748}
      \strng{authorfullhash}{5438078037fd2ae6b07b804b456fc748}
      \field{extraname}{1}
      \field{labelalpha}{BS22}
      \field{sortinit}{B}
      \field{sortinithash}{276475738cc058478c1677046f857703}
      \field{labelnamesource}{author}
      \field{labeltitlesource}{title}
      \field{abstract}{We study the solutions of the one-phase supercooled Stefan problem with kinetic undercooling, which describes the freezing of a supercooled liquid, in one spatial dimension. Assuming that the initial temperature lies between the equilibrium freezing point and the characteristic invariant temperature throughout the liquid our main theorem shows that, as the kinetic undercooling parameter tends to zero, the free boundary converges to the (possibly irregular) free boundary in the supercooled Stefan problem without kinetic undercooling, whose uniqueness has been recently established in (Delarue, Nadtochiy and Shkolnikov (2019), Ledger and Søjmark (2018)). The key tools in the proof are a Feynman–Kac formula, which expresses the free boundary in the problem with kinetic undercooling through a local time of a reflected process, and a resulting comparison principle for the free boundaries with different kinetic undercooling parameters.}
      \field{issn}{0246-0203}
      \field{journaltitle}{Annales de l'Institut Henri Poincaré, Probabilités et Statistiques}
      \field{month}{5}
      \field{number}{2}
      \field{title}{Zero kinetic undercooling limit in the supercooled {Stefan} problem}
      \field{urlday}{4}
      \field{urlmonth}{11}
      \field{urlyear}{2022}
      \field{volume}{58}
      \field{year}{2022}
      \field{urldateera}{ce}
      \verb{doi}
      \verb 10.1214/21-AIHP1194
      \endverb
      \verb{file}
      \verb Baker and Shkolnikov - 2022 - Zero kinetic undercooling limit in the supercooled.pdf:/Users/graemeb/Zotero/storage/FN3KLAEU/Baker and Shkolnikov - 2022 - Zero kinetic undercooling limit in the supercooled.pdf:application/pdf
      \endverb
      \verb{urlraw}
      \verb https://projecteuclid.org/journals/annales-de-linstitut-henri-poincare-probabilites-et-statistiques/volume-58/issue-2/Zero-kinetic-undercooling-limit-in-the-supercooled-Stefan-problem/10.1214/21-AIHP1194.full
      \endverb
      \verb{url}
      \verb https://projecteuclid.org/journals/annales-de-linstitut-henri-poincare-probabilites-et-statistiques/volume-58/issue-2/Zero-kinetic-undercooling-limit-in-the-supercooled-Stefan-problem/10.1214/21-AIHP1194.full
      \endverb
    \endentry
    \entry{baker_singular_2023}{misc}{}
      \name{author}{2}{}{%
        {{hash=0761dfc2a74ec6f0d1effbf27f715d59}{%
           family={Baker},
           familyi={B\bibinitperiod},
           given={Graeme},
           giveni={G\bibinitperiod}}}%
        {{hash=cce28fd978c3edebbc2aee9c6f7e9b0f}{%
           family={Shkolnikov},
           familyi={S\bibinitperiod},
           given={Mykhaylo},
           giveni={M\bibinitperiod}}}%
      }
      \list{publisher}{1}{%
        {arXiv}%
      }
      \strng{namehash}{5438078037fd2ae6b07b804b456fc748}
      \strng{fullhash}{5438078037fd2ae6b07b804b456fc748}
      \strng{bibnamehash}{5438078037fd2ae6b07b804b456fc748}
      \strng{authorbibnamehash}{5438078037fd2ae6b07b804b456fc748}
      \strng{authornamehash}{5438078037fd2ae6b07b804b456fc748}
      \strng{authorfullhash}{5438078037fd2ae6b07b804b456fc748}
      \field{extraname}{2}
      \field{labelalpha}{BS23}
      \field{sortinit}{B}
      \field{sortinithash}{276475738cc058478c1677046f857703}
      \field{labelnamesource}{author}
      \field{labeltitlesource}{title}
      \field{abstract}{We consider a probabilistic formulation of a singular two-phase Stefan problem in one space dimension, which amounts to a coupled system of two McKean-Vlasov stochastic differential equations. In the financial context of systemic risk, this system models two competing regions with a large number of interconnected banks or firms at risk of default. Our main result shows the existence of a solution whose discontinuities obey the natural physicality condition for the problem at hand. Thus, this work extends the recent series of existence results for singular one-phase Stefan problems in one space dimension that can be found in [DIRT15a], [NS19a], [HLS18], [CRS20]. As therein, our existence result is obtained via a large system limit of a finite particle system approximation in the Skorokhod M1 topology. But, unlike for the previously studied one-phase case, the free boundary herein is not monotone, so that the large system limit is obtained by a novel argument.}
      \field{month}{4}
      \field{note}{arXiv:2203.06003 [math-ph]}
      \field{title}{A singular two-phase {Stefan} problem and particles interacting through their hitting times}
      \field{urlday}{5}
      \field{urlmonth}{5}
      \field{urlyear}{2023}
      \field{year}{2023}
      \field{urldateera}{ce}
      \verb{doi}
      \verb 10.48550/arXiv.2203.06003
      \endverb
      \verb{file}
      \verb arXiv Fulltext PDF:/Users/graemeb/Zotero/storage/9XI66U2Q/Baker and Shkolnikov - 2023 - A singular two-phase Stefan problem and particles .pdf:application/pdf;arXiv.org Snapshot:/Users/graemeb/Zotero/storage/7FE6Z2DC/2203.html:text/html
      \endverb
      \verb{urlraw}
      \verb http://arxiv.org/abs/2203.06003
      \endverb
      \verb{url}
      \verb http://arxiv.org/abs/2203.06003
      \endverb
      \keyw{82C22,60H30,35K10,35K60,91G80,Mathematical Physics,Mathematics - Analysis of PDEs,Mathematics - Probability}
    \endentry
    \entry{cuchiero_propagation_2020}{misc}{}
      \name{author}{3}{}{%
        {{hash=64fb2fa55e95db41a74d9e72b25c918e}{%
           family={Cuchiero},
           familyi={C\bibinitperiod},
           given={Christa},
           giveni={C\bibinitperiod}}}%
        {{hash=1d6ef4334530eed497fe0cc823ce5ef0}{%
           family={Rigger},
           familyi={R\bibinitperiod},
           given={Stefan},
           giveni={S\bibinitperiod}}}%
        {{hash=d624c4119e82128ac097fec1bef6c379}{%
           family={Svaluto-Ferro},
           familyi={S\bibinithyphendelim F\bibinitperiod},
           given={Sara},
           giveni={S\bibinitperiod}}}%
      }
      \list{publisher}{1}{%
        {arXiv}%
      }
      \strng{namehash}{cd15d99697153a3cb25b1cc45def44a6}
      \strng{fullhash}{cd15d99697153a3cb25b1cc45def44a6}
      \strng{bibnamehash}{cd15d99697153a3cb25b1cc45def44a6}
      \strng{authorbibnamehash}{cd15d99697153a3cb25b1cc45def44a6}
      \strng{authornamehash}{cd15d99697153a3cb25b1cc45def44a6}
      \strng{authorfullhash}{cd15d99697153a3cb25b1cc45def44a6}
      \field{labelalpha}{CRS20}
      \field{sortinit}{C}
      \field{sortinithash}{963e9d84a3da2344e8833203de5aed05}
      \field{labelnamesource}{author}
      \field{labeltitlesource}{title}
      \field{abstract}{Supercooled Stefan problems describe the evolution of the boundary between the solid and liquid phases of a substance, where the liquid is assumed to be cooled below its freezing point. Following the methodology of Delarue, Nadtochiy and Shkolnikov, we construct solutions to the one-phase one-dimensional supercooled Stefan problem through a certain McKean-Vlasov equation, which allows to define global solutions even in the presence of blow-ups. Solutions to the McKean-Vlasov equation arise as mean-field limits of particle systems interacting through hitting times, which is important for systemic risk modeling. Our main contributions are: (i) we prove a general tightness theorem for the Skorokhod M1-topology which applies to processes that can be decomposed into a continuous and a monotone part. (ii) We prove propagation of chaos for a perturbed version of the particle system for general initial conditions. (iii) We prove a conjecture of Delarue, Nadtochiy and Shkolnikov, relating the solution concepts of so-called minimal and physical solutions, showing that minimal solutions of the McKean-Vlasov equation are physical whenever the initial condition is integrable.}
      \field{month}{10}
      \field{note}{arXiv:2010.03580 [math, q-fin] version: 1}
      \field{title}{Propagation of minimality in the supercooled {Stefan} problem}
      \field{urlday}{20}
      \field{urlmonth}{6}
      \field{urlyear}{2023}
      \field{year}{2020}
      \field{urldateera}{ce}
      \verb{doi}
      \verb 10.48550/arXiv.2010.03580
      \endverb
      \verb{file}
      \verb arXiv Fulltext PDF:/Users/graemeb/Zotero/storage/E9TFVEC6/Cuchiero et al. - 2020 - Propagation of minimality in the supercooled Stefa.pdf:application/pdf;arXiv.org Snapshot:/Users/graemeb/Zotero/storage/TSFL88UC/2010.html:text/html
      \endverb
      \verb{urlraw}
      \verb http://arxiv.org/abs/2010.03580
      \endverb
      \verb{url}
      \verb http://arxiv.org/abs/2010.03580
      \endverb
      \keyw{Mathematics - Probability,Quantitative Finance - Mathematical Finance}
    \endentry
    \entry{delarue_particle_2015}{misc}{}
      \name{author}{4}{}{%
        {{hash=4f0c06214bb0b35a123f8990c42412bc}{%
           family={Delarue},
           familyi={D\bibinitperiod},
           given={F.},
           giveni={F\bibinitperiod}}}%
        {{hash=b7234b63caa4c77e6969cf97105c3359}{%
           family={Inglis},
           familyi={I\bibinitperiod},
           given={J.},
           giveni={J\bibinitperiod}}}%
        {{hash=3b5a9f87943f0cf3dc9085bd9fd4de7d}{%
           family={Rubenthaler},
           familyi={R\bibinitperiod},
           given={S.},
           giveni={S\bibinitperiod}}}%
        {{hash=6c29d4f0daf02fc9c1e7a7a49e652d2e}{%
           family={Tanré},
           familyi={T\bibinitperiod},
           given={E.},
           giveni={E\bibinitperiod}}}%
      }
      \list{language}{1}{%
        {en}%
      }
      \strng{namehash}{5d4a6862aeeba873986dd19d5a94299b}
      \strng{fullhash}{5d4a6862aeeba873986dd19d5a94299b}
      \strng{bibnamehash}{5d4a6862aeeba873986dd19d5a94299b}
      \strng{authorbibnamehash}{5d4a6862aeeba873986dd19d5a94299b}
      \strng{authornamehash}{5d4a6862aeeba873986dd19d5a94299b}
      \strng{authorfullhash}{5d4a6862aeeba873986dd19d5a94299b}
      \field{labelalpha}{DIRT15}
      \field{sortinit}{D}
      \field{sortinithash}{2ef1bd9a78cc71eb74d7231c635177b8}
      \field{labelnamesource}{author}
      \field{labeltitlesource}{title}
      \field{abstract}{We discuss the construction and approximation of solutions to a nonlinear McKean-Vlasov equation driven by a singular self-excitatory interaction of the meanfield type. Such an equation is intended to describe an infinite population of neurons which interact with one another. Each time a proportion of neurons ‘spike’, the whole network instantaneously receives an excitatory kick. The instantaneous nature of the excitation makes the system singular and prevents the application of standard results from the literature. Making use of the Skorohod M1 topology, we prove that, for the right notion of a ‘physical’ solution, the nonlinear equation can be approximated either by a finite particle system or by a delayed equation. As a by-product, we obtain the existence of ‘synchronized’ solutions, for which a macroscopic proportion of neurons may spike at the same time.}
      \field{month}{1}
      \field{note}{arXiv: 1406.1151}
      \field{title}{Particle systems with a singular mean-field self-excitation. {Application} to neuronal networks}
      \field{urlday}{3}
      \field{urlmonth}{8}
      \field{urlyear}{2020}
      \field{year}{2015}
      \field{urldateera}{ce}
      \verb{file}
      \verb Delarue et al. - 2015 - Particle systems with a singular mean-field self-e.pdf:/Users/graemeb/Zotero/storage/3B8383JM/Delarue et al. - 2015 - Particle systems with a singular mean-field self-e.pdf:application/pdf
      \endverb
      \verb{urlraw}
      \verb http://arxiv.org/abs/1406.1151
      \endverb
      \verb{url}
      \verb http://arxiv.org/abs/1406.1151
      \endverb
      \keyw{Mathematics - Probability}
    \endentry
    \entry{delarue_global_2022}{article}{}
      \name{author}{3}{}{%
        {{hash=799c6f3ae3f94478b9b2d2aadabd8387}{%
           family={Delarue},
           familyi={D\bibinitperiod},
           given={François},
           giveni={F\bibinitperiod}}}%
        {{hash=23bde0d98049dfa33360fa0880afaec5}{%
           family={Nadtochiy},
           familyi={N\bibinitperiod},
           given={Sergey},
           giveni={S\bibinitperiod}}}%
        {{hash=cce28fd978c3edebbc2aee9c6f7e9b0f}{%
           family={Shkolnikov},
           familyi={S\bibinitperiod},
           given={Mykhaylo},
           giveni={M\bibinitperiod}}}%
      }
      \list{language}{1}{%
        {en}%
      }
      \strng{namehash}{7d3c61a499f82754336eede9cec884c7}
      \strng{fullhash}{7d3c61a499f82754336eede9cec884c7}
      \strng{bibnamehash}{7d3c61a499f82754336eede9cec884c7}
      \strng{authorbibnamehash}{7d3c61a499f82754336eede9cec884c7}
      \strng{authornamehash}{7d3c61a499f82754336eede9cec884c7}
      \strng{authorfullhash}{7d3c61a499f82754336eede9cec884c7}
      \field{labelalpha}{DNS22}
      \field{sortinit}{D}
      \field{sortinithash}{2ef1bd9a78cc71eb74d7231c635177b8}
      \field{labelnamesource}{author}
      \field{labeltitlesource}{shorttitle}
      \field{issn}{2690-1005, 2690-0998}
      \field{journaltitle}{Probability and Mathematical Physics}
      \field{month}{5}
      \field{number}{1}
      \field{shorttitle}{Global solutions to the supercooled {Stefan} problem with blow-ups}
      \field{title}{Global solutions to the supercooled {Stefan} problem with blow-ups: regularity and uniqueness}
      \field{urlday}{3}
      \field{urlmonth}{1}
      \field{urlyear}{2023}
      \field{volume}{3}
      \field{year}{2022}
      \field{urldateera}{ce}
      \field{pages}{171\bibrangedash 213}
      \range{pages}{43}
      \verb{doi}
      \verb 10.2140/pmp.2022.3.171
      \endverb
      \verb{file}
      \verb Delarue et al. - 2022 - Global solutions to the supercooled Stefan problem.pdf:/Users/graemeb/Zotero/storage/MKV87BJY/Delarue et al. - 2022 - Global solutions to the supercooled Stefan problem.pdf:application/pdf
      \endverb
      \verb{urlraw}
      \verb https://msp.org/pmp/2022/3-1/p06.xhtml
      \endverb
      \verb{url}
      \verb https://msp.org/pmp/2022/3-1/p06.xhtml
      \endverb
    \endentry
    \entry{dudley_real_2002}{book}{}
      \name{author}{1}{}{%
        {{hash=fbefc6b548928472eeb7c87d77ea10cf}{%
           family={Dudley},
           familyi={D\bibinitperiod},
           given={Richard\bibnamedelima M},
           giveni={R\bibinitperiod\bibinitdelim M\bibinitperiod}}}%
      }
      \list{publisher}{1}{%
        {Cambridge University Press}%
      }
      \strng{namehash}{fbefc6b548928472eeb7c87d77ea10cf}
      \strng{fullhash}{fbefc6b548928472eeb7c87d77ea10cf}
      \strng{bibnamehash}{fbefc6b548928472eeb7c87d77ea10cf}
      \strng{authorbibnamehash}{fbefc6b548928472eeb7c87d77ea10cf}
      \strng{authornamehash}{fbefc6b548928472eeb7c87d77ea10cf}
      \strng{authorfullhash}{fbefc6b548928472eeb7c87d77ea10cf}
      \field{labelalpha}{Dud02}
      \field{sortinit}{D}
      \field{sortinithash}{2ef1bd9a78cc71eb74d7231c635177b8}
      \field{labelnamesource}{author}
      \field{labeltitlesource}{title}
      \field{title}{Real analysis and probability}
      \field{year}{2002}
    \endentry
    \entry{evans_partial_2010}{book}{}
      \name{author}{1}{}{%
        {{hash=0d72932a7775f85a51622d3af75ae26e}{%
           family={Evans},
           familyi={E\bibinitperiod},
           given={Lawrence\bibnamedelima C},
           giveni={L\bibinitperiod\bibinitdelim C\bibinitperiod}}}%
      }
      \list{publisher}{1}{%
        {American Mathematical Socety}%
      }
      \strng{namehash}{0d72932a7775f85a51622d3af75ae26e}
      \strng{fullhash}{0d72932a7775f85a51622d3af75ae26e}
      \strng{bibnamehash}{0d72932a7775f85a51622d3af75ae26e}
      \strng{authorbibnamehash}{0d72932a7775f85a51622d3af75ae26e}
      \strng{authornamehash}{0d72932a7775f85a51622d3af75ae26e}
      \strng{authorfullhash}{0d72932a7775f85a51622d3af75ae26e}
      \field{labelalpha}{Eva10}
      \field{sortinit}{E}
      \field{sortinithash}{f615fb9c6fba11c6f962fb3fd599810e}
      \field{labelnamesource}{author}
      \field{labeltitlesource}{title}
      \field{title}{Partial {Differential} {Equations}}
      \field{volume}{19}
      \field{year}{2010}
    \endentry
    \entry{hambly_mckean--vlasov_2018}{misc}{}
      \name{author}{3}{}{%
        {{hash=6d9dd5545b1baec4468e75511074ebd7}{%
           family={Hambly},
           familyi={H\bibinitperiod},
           given={Ben},
           giveni={B\bibinitperiod}}}%
        {{hash=fd058415ea8182ef4b437593a3f6b573}{%
           family={Ledger},
           familyi={L\bibinitperiod},
           given={Sean},
           giveni={S\bibinitperiod}}}%
        {{hash=af7642936d68b11531642ba1fc8e78de}{%
           family={Sojmark},
           familyi={S\bibinitperiod},
           given={Andreas},
           giveni={A\bibinitperiod}}}%
      }
      \list{publisher}{1}{%
        {arXiv}%
      }
      \strng{namehash}{86ea65ad0ddb64b18cb0019da1d0bea7}
      \strng{fullhash}{86ea65ad0ddb64b18cb0019da1d0bea7}
      \strng{bibnamehash}{86ea65ad0ddb64b18cb0019da1d0bea7}
      \strng{authorbibnamehash}{86ea65ad0ddb64b18cb0019da1d0bea7}
      \strng{authornamehash}{86ea65ad0ddb64b18cb0019da1d0bea7}
      \strng{authorfullhash}{86ea65ad0ddb64b18cb0019da1d0bea7}
      \field{labelalpha}{HLS18}
      \field{sortinit}{H}
      \field{sortinithash}{5f15a7bc777ad49ff15aa4d2831b1681}
      \field{labelnamesource}{author}
      \field{labeltitlesource}{title}
      \field{abstract}{We study a McKean--Vlasov equation arising from a mean-field model of a particle system with positive feedback. As particles hit a barrier they cause the other particles to jump in the direction of the barrier and this feedback mechanism leads to the possibility that the system can exhibit contagious blow-ups. Using a fixed-point argument we construct a differentiable solution up to a first explosion time. Our main contribution is a proof of uniqueness in the class of c{\textbackslash}`\{a\}dl{\textbackslash}`\{a\}g functions, which confirms the validity of related propagation-of-chaos results in the literature. We extend the allowed initial conditions to include densities with any power law decay at the boundary, and connect the exponent of decay with the growth exponent of the solution in small time in a precise way. This takes us asymptotically close to the control on initial conditions required for a global solution theory. A novel minimality result and trapping technique are introduced to prove uniqueness.}
      \field{month}{10}
      \field{note}{arXiv:1801.07703 [math]}
      \field{title}{A {McKean}--{Vlasov} equation with positive feedback and blow-ups}
      \field{urlday}{20}
      \field{urlmonth}{6}
      \field{urlyear}{2023}
      \field{year}{2018}
      \field{urldateera}{ce}
      \verb{doi}
      \verb 10.48550/arXiv.1801.07703
      \endverb
      \verb{file}
      \verb arXiv Fulltext PDF:/Users/graemeb/Zotero/storage/WZHLX4VA/Hambly et al. - 2018 - A McKean--Vlasov equation with positive feedback a.pdf:application/pdf;arXiv.org Snapshot:/Users/graemeb/Zotero/storage/KZLHPP9X/1801.html:text/html
      \endverb
      \verb{urlraw}
      \verb http://arxiv.org/abs/1801.07703
      \endverb
      \verb{url}
      \verb http://arxiv.org/abs/1801.07703
      \endverb
      \keyw{60H10,60H30,Mathematics - Analysis of PDEs,Mathematics - Probability}
    \endentry
    \entry{hambly_mckean-vlasov_2022}{misc}{}
      \name{author}{2}{}{%
        {{hash=6d9dd5545b1baec4468e75511074ebd7}{%
           family={Hambly},
           familyi={H\bibinitperiod},
           given={Ben},
           giveni={B\bibinitperiod}}}%
        {{hash=8a35cb13933201cd49960d09ff7f252d}{%
           family={Meier},
           familyi={M\bibinitperiod},
           given={Julian},
           giveni={J\bibinitperiod}}}%
      }
      \list{publisher}{1}{%
        {arXiv}%
      }
      \strng{namehash}{cfb0839e857c5d5566f53e578aaf8a5b}
      \strng{fullhash}{cfb0839e857c5d5566f53e578aaf8a5b}
      \strng{bibnamehash}{cfb0839e857c5d5566f53e578aaf8a5b}
      \strng{authorbibnamehash}{cfb0839e857c5d5566f53e578aaf8a5b}
      \strng{authornamehash}{cfb0839e857c5d5566f53e578aaf8a5b}
      \strng{authorfullhash}{cfb0839e857c5d5566f53e578aaf8a5b}
      \field{labelalpha}{HM22}
      \field{sortinit}{H}
      \field{sortinithash}{5f15a7bc777ad49ff15aa4d2831b1681}
      \field{labelnamesource}{author}
      \field{labeltitlesource}{title}
      \field{abstract}{We prove existence and uniqueness of physical and minimal solutions to McKean-Vlasov equations with positive feedback through elastic stopping times. We do this by establishing a relationship between this problem and a problem with absorbing stopping times. We show convergence of a particle system to the McKean-Vlasov equation. Moreover, we establish convergence of the elastic McKean-Vlasov problem to the problem with absorbing stopping times and to a reflecting Brownian motion as the elastic parameter goes to infinity or zero respectively.}
      \field{month}{3}
      \field{note}{arXiv:2203.10787 [math]}
      \field{title}{{McKean}-{Vlasov} {Equations} with {Positive} {Feedback} through {Elastic} {Stopping} {Times}}
      \field{urlday}{20}
      \field{urlmonth}{6}
      \field{urlyear}{2023}
      \field{year}{2022}
      \field{urldateera}{ce}
      \verb{doi}
      \verb 10.48550/arXiv.2203.10787
      \endverb
      \verb{file}
      \verb arXiv Fulltext PDF:/Users/graemeb/Zotero/storage/TV6HHDQJ/Hambly and Meier - 2022 - McKean-Vlasov Equations with Positive Feedback thr.pdf:application/pdf;arXiv.org Snapshot:/Users/graemeb/Zotero/storage/N9HW8XTU/2203.html:text/html
      \endverb
      \verb{urlraw}
      \verb http://arxiv.org/abs/2203.10787
      \endverb
      \verb{url}
      \verb http://arxiv.org/abs/2203.10787
      \endverb
      \keyw{35K61,60H30,60K35,Mathematics - Probability}
    \endentry
    \entry{hambly_spde_2019}{article}{}
      \name{author}{2}{}{%
        {{hash=6d9dd5545b1baec4468e75511074ebd7}{%
           family={Hambly},
           familyi={H\bibinitperiod},
           given={Ben},
           giveni={B\bibinitperiod}}}%
        {{hash=af6596c726b8d0dca3dbf8d39a9fc0d7}{%
           family={Søjmark},
           familyi={S\bibinitperiod},
           given={Andreas},
           giveni={A\bibinitperiod}}}%
      }
      \list{language}{1}{%
        {en}%
      }
      \strng{namehash}{82691804109eca2ff6cc6579163cff11}
      \strng{fullhash}{82691804109eca2ff6cc6579163cff11}
      \strng{bibnamehash}{82691804109eca2ff6cc6579163cff11}
      \strng{authorbibnamehash}{82691804109eca2ff6cc6579163cff11}
      \strng{authornamehash}{82691804109eca2ff6cc6579163cff11}
      \strng{authorfullhash}{82691804109eca2ff6cc6579163cff11}
      \field{labelalpha}{HS19}
      \field{sortinit}{H}
      \field{sortinithash}{5f15a7bc777ad49ff15aa4d2831b1681}
      \field{labelnamesource}{author}
      \field{labeltitlesource}{title}
      \field{abstract}{We propose a dynamic mean-field model for ‘systemic risk’ in large financial systems, derived from a system of interacting diffusions on the positive half-line with an absorbing boundary at the origin. These diffusions represent the distancesto-default of financial institutions, and absorption at zero corresponds to default. As a way of modelling correlated exposures and herd behaviour, we consider a common source of noise and a form of mean-reversion in the drift. Moreover, we introduce an endogenous contagion mechanism whereby the default of one institution causes a drop in the distances-to-default of the other institutions. In this way, we aim to capture key ‘system-wide’ effects on risk. The resulting mean-field limit is characterised uniquely by a nonlinear SPDE on the half-line with a Dirichlet boundary condition. The density of this SPDE gives the conditional law of a non-standard ‘conditional’ McKean–Vlasov diffusion, for which we provide a novel upper Dirichlet heat kernel type estimate. Depending on the realisations of the common noise and the rate of mean-reversion, the SPDE can exhibit rapid accelerations in the loss of mass at the boundary. In other words, the contagion mechanism can give rise to periods of significant systemic default clustering.}
      \field{issn}{0949-2984, 1432-1122}
      \field{journaltitle}{Finance and Stochastics}
      \field{month}{7}
      \field{number}{3}
      \field{title}{An {SPDE} model for systemic risk with endogenous contagion}
      \field{urlday}{21}
      \field{urlmonth}{11}
      \field{urlyear}{2022}
      \field{volume}{23}
      \field{year}{2019}
      \field{urldateera}{ce}
      \field{pages}{535\bibrangedash 594}
      \range{pages}{60}
      \verb{doi}
      \verb 10.1007/s00780-019-00396-1
      \endverb
      \verb{file}
      \verb Hambly and Søjmark - 2019 - An SPDE model for systemic risk with endogenous co.pdf:/Users/graemeb/Zotero/storage/64CX3YQ9/Hambly and Søjmark - 2019 - An SPDE model for systemic risk with endogenous co.pdf:application/pdf
      \endverb
      \verb{urlraw}
      \verb http://link.springer.com/10.1007/s00780-019-00396-1
      \endverb
      \verb{url}
      \verb http://link.springer.com/10.1007/s00780-019-00396-1
      \endverb
    \endentry
    \entry{ledger_uniqueness_2020}{article}{}
      \name{author}{2}{}{%
        {{hash=fd058415ea8182ef4b437593a3f6b573}{%
           family={Ledger},
           familyi={L\bibinitperiod},
           given={Sean},
           giveni={S\bibinitperiod}}}%
        {{hash=af6596c726b8d0dca3dbf8d39a9fc0d7}{%
           family={Søjmark},
           familyi={S\bibinitperiod},
           given={Andreas},
           giveni={A\bibinitperiod}}}%
      }
      \list{language}{1}{%
        {en}%
      }
      \strng{namehash}{44b23ced1f0e543328e8801a084cff3c}
      \strng{fullhash}{44b23ced1f0e543328e8801a084cff3c}
      \strng{bibnamehash}{44b23ced1f0e543328e8801a084cff3c}
      \strng{authorbibnamehash}{44b23ced1f0e543328e8801a084cff3c}
      \strng{authornamehash}{44b23ced1f0e543328e8801a084cff3c}
      \strng{authorfullhash}{44b23ced1f0e543328e8801a084cff3c}
      \field{labelalpha}{LS20}
      \field{sortinit}{L}
      \field{sortinithash}{2c7981aaabc885868aba60f0c09ee20f}
      \field{labelnamesource}{author}
      \field{labeltitlesource}{title}
      \field{abstract}{We present a simple uniqueness argument for a collection of McKean–Vlasov problems that have seen recent interest. Our first result shows that, in the weak feedback regime, there is global uniqueness for a very general class of random drivers. By weak feedback we mean the case where the contagion parameters are small enough to prevent blow-ups in solutions. Next, we specialise to a Brownian driver and show how the same techniques can be extended to give short-time uniqueness after blow-ups, regardless of the feedback strength, by exploiting a localised form of weak feedback. The heart of our approach is a surprisingly simple probabilistic comparison argument that is robust in the sense that it does not ask for any regularity of the solutions.}
      \field{issn}{0024-6093, 1469-2120}
      \field{journaltitle}{Bulletin of the London Mathematical Society}
      \field{month}{6}
      \field{number}{3}
      \field{title}{Uniqueness for contagious {McKean}–{Vlasov} systems in the weak feedback regime}
      \field{urlday}{29}
      \field{urlmonth}{7}
      \field{urlyear}{2021}
      \field{volume}{52}
      \field{year}{2020}
      \field{urldateera}{ce}
      \field{pages}{448\bibrangedash 463}
      \range{pages}{16}
      \verb{doi}
      \verb 10.1112/blms.12337
      \endverb
      \verb{file}
      \verb Ledger and Søjmark - 2020 - Uniqueness for contagious McKean–Vlasov systems in.pdf:/Users/graemeb/Zotero/storage/8PR6C7AS/Ledger and Søjmark - 2020 - Uniqueness for contagious McKean–Vlasov systems in.pdf:application/pdf
      \endverb
      \verb{urlraw}
      \verb https://onlinelibrary.wiley.com/doi/10.1112/blms.12337
      \endverb
      \verb{url}
      \verb https://onlinelibrary.wiley.com/doi/10.1112/blms.12337
      \endverb
    \endentry
    \entry{mustapha_well-posedness_2023}{misc}{}
      \name{author}{2}{}{%
        {{hash=472df37ad5ebe2c672717e5df8f26e4d}{%
           family={Mustapha},
           familyi={M\bibinitperiod},
           given={Scander},
           giveni={S\bibinitperiod}}}%
        {{hash=cce28fd978c3edebbc2aee9c6f7e9b0f}{%
           family={Shkolnikov},
           familyi={S\bibinitperiod},
           given={Mykhaylo},
           giveni={M\bibinitperiod}}}%
      }
      \list{publisher}{1}{%
        {arXiv}%
      }
      \strng{namehash}{8caa8fbaffad08302664fcadd7da619f}
      \strng{fullhash}{8caa8fbaffad08302664fcadd7da619f}
      \strng{bibnamehash}{8caa8fbaffad08302664fcadd7da619f}
      \strng{authorbibnamehash}{8caa8fbaffad08302664fcadd7da619f}
      \strng{authornamehash}{8caa8fbaffad08302664fcadd7da619f}
      \strng{authorfullhash}{8caa8fbaffad08302664fcadd7da619f}
      \field{labelalpha}{MS23}
      \field{sortinit}{M}
      \field{sortinithash}{cfd219b90152c06204fab207bc6c7cab}
      \field{labelnamesource}{author}
      \field{labeltitlesource}{title}
      \field{abstract}{We study the one-phase one-dimensional supercooled Stefan problem with oscillatory initial conditions. In this context, the global existence of so-called physical solutions has been shown recently in [CRSF20], despite the presence of blow-ups in the freezing rate. On the other hand, for regular initial conditions, the uniqueness of physical solutions has been established in [DNS22]. Here, we prove the uniqueness of physical solutions for oscillatory initial conditions by a new contraction argument that allows to replace the local monotonicity condition of [DNS22] with an averaging condition. We verify this weaker condition for fairly general oscillating probability densities, such as the ones given by an almost sure trajectory of \$(1+W\_x-{\textbackslash}sqrt\{2x{|}{\textbackslash}log\{{|}{\textbackslash}log\{x\}{|}\}{|}\})\_\{+\}{\textbackslash}wedge 1\$ near the origin, where \$W\$ is a standard Brownian motion. We also permit typical deterministically constructed oscillating densities, including those of the form \$(1+{\textbackslash}sin\{1/x\})/2\$ near the origin. Finally, we provide an example of oscillating densities for which it is possible to go beyond our main assumption via further complementary arguments.}
      \field{month}{4}
      \field{note}{arXiv:2302.13097 [math-ph]}
      \field{title}{Well-posedness of the supercooled {Stefan} problem with oscillatory initial conditions}
      \field{urlday}{20}
      \field{urlmonth}{6}
      \field{urlyear}{2023}
      \field{year}{2023}
      \field{urldateera}{ce}
      \verb{doi}
      \verb 10.48550/arXiv.2302.13097
      \endverb
      \verb{file}
      \verb arXiv Fulltext PDF:/Users/graemeb/Zotero/storage/NMEW3GMF/Mustapha and Shkolnikov - 2023 - Well-posedness of the supercooled Stefan problem w.pdf:application/pdf;arXiv.org Snapshot:/Users/graemeb/Zotero/storage/TREM9FY6/2302.html:text/html
      \endverb
      \verb{urlraw}
      \verb http://arxiv.org/abs/2302.13097
      \endverb
      \verb{url}
      \verb http://arxiv.org/abs/2302.13097
      \endverb
      \keyw{80A22,35B05,60H30,Mathematical Physics,Mathematics - Analysis of PDEs,Mathematics - Probability}
    \endentry
    \entry{nadtochiy_mean_2019}{misc}{}
      \name{author}{2}{}{%
        {{hash=23bde0d98049dfa33360fa0880afaec5}{%
           family={Nadtochiy},
           familyi={N\bibinitperiod},
           given={Sergey},
           giveni={S\bibinitperiod}}}%
        {{hash=cce28fd978c3edebbc2aee9c6f7e9b0f}{%
           family={Shkolnikov},
           familyi={S\bibinitperiod},
           given={Mykhaylo},
           giveni={M\bibinitperiod}}}%
      }
      \list{publisher}{1}{%
        {arXiv}%
      }
      \strng{namehash}{a0590b36e4b9977c6746f946d087f6c1}
      \strng{fullhash}{a0590b36e4b9977c6746f946d087f6c1}
      \strng{bibnamehash}{a0590b36e4b9977c6746f946d087f6c1}
      \strng{authorbibnamehash}{a0590b36e4b9977c6746f946d087f6c1}
      \strng{authornamehash}{a0590b36e4b9977c6746f946d087f6c1}
      \strng{authorfullhash}{a0590b36e4b9977c6746f946d087f6c1}
      \field{extraname}{1}
      \field{labelalpha}{NS19}
      \field{sortinit}{N}
      \field{sortinithash}{f7242c3ed3dc50029fca1be76c497c7c}
      \field{extraalpha}{1}
      \field{labelnamesource}{author}
      \field{labeltitlesource}{title}
      \field{abstract}{Building on the line of work [DIRT15a], [DIRT15b], [NS17a], [DT17], [HLS18], [HS18] we continue the study of particle systems with singular interaction through hitting times. In contrast to the previous research, we (i) consider very general driving processes and interaction functions, (ii) allow for inhomogeneous connection structures, and (iii) analyze a game in which the particles determine their connections strategically. Hereby, we uncover two completely new phenomena. First, we characterize the "times of fragility" of such systems (e.g., the times when a macroscopic part of the population defaults or gets infected simultaneously, or when the neuron cells "synchronize") explicitly in terms of the dynamics of the driving processes, the current distribution of the particles' values, and the topology of the underlying network (represented by its Perron-Frobenius eigenvalue). Second, we use such systems to describe a dynamic credit-network game and show that, in equilibrium, the system regularizes: i.e., the times of fragility never occur, as the particles avoid them by adjusting their connections strategically. Two auxiliary mathematical results, useful in their own right, are uncovered during our investigation: a generalization of Schauder's fixed-point theorem for the Skorokhod space with the M1 topology, and the application of the max-plus algebra to the equilibrium version of the network flow problem.}
      \field{month}{9}
      \field{note}{arXiv:1807.02015 [math, q-fin]}
      \field{title}{Mean field systems on networks, with singular interaction through hitting times}
      \field{urlday}{20}
      \field{urlmonth}{6}
      \field{urlyear}{2023}
      \field{year}{2019}
      \field{urldateera}{ce}
      \verb{doi}
      \verb 10.48550/arXiv.1807.02015
      \endverb
      \verb{file}
      \verb arXiv Fulltext PDF:/Users/graemeb/Zotero/storage/NT8EILLD/Nadtochiy and Shkolnikov - 2019 - Mean field systems on networks, with singular inte.pdf:application/pdf;arXiv.org Snapshot:/Users/graemeb/Zotero/storage/GS2HHATA/1807.html:text/html
      \endverb
      \verb{urlraw}
      \verb http://arxiv.org/abs/1807.02015
      \endverb
      \verb{url}
      \verb http://arxiv.org/abs/1807.02015
      \endverb
      \keyw{82C22,91A25,35K40,54H25,15A80,Mathematics - Analysis of PDEs,Mathematics - General Topology,Mathematics - Probability,Quantitative Finance - Mathematical Finance}
    \endentry
    \entry{nadtochiy_particle_2019}{article}{}
      \name{author}{2}{}{%
        {{hash=23bde0d98049dfa33360fa0880afaec5}{%
           family={Nadtochiy},
           familyi={N\bibinitperiod},
           given={Sergey},
           giveni={S\bibinitperiod}}}%
        {{hash=cce28fd978c3edebbc2aee9c6f7e9b0f}{%
           family={Shkolnikov},
           familyi={S\bibinitperiod},
           given={Mykhaylo},
           giveni={M\bibinitperiod}}}%
      }
      \list{language}{1}{%
        {en}%
      }
      \strng{namehash}{a0590b36e4b9977c6746f946d087f6c1}
      \strng{fullhash}{a0590b36e4b9977c6746f946d087f6c1}
      \strng{bibnamehash}{a0590b36e4b9977c6746f946d087f6c1}
      \strng{authorbibnamehash}{a0590b36e4b9977c6746f946d087f6c1}
      \strng{authornamehash}{a0590b36e4b9977c6746f946d087f6c1}
      \strng{authorfullhash}{a0590b36e4b9977c6746f946d087f6c1}
      \field{extraname}{2}
      \field{labelalpha}{NS19}
      \field{sortinit}{N}
      \field{sortinithash}{f7242c3ed3dc50029fca1be76c497c7c}
      \field{extraalpha}{2}
      \field{labelnamesource}{author}
      \field{labeltitlesource}{shorttitle}
      \field{issn}{1050-5164}
      \field{journaltitle}{The Annals of Applied Probability}
      \field{month}{2}
      \field{number}{1}
      \field{shorttitle}{Particle systems with singular interaction through hitting times}
      \field{title}{Particle systems with singular interaction through hitting times: {Application} in systemic risk modeling}
      \field{urlday}{3}
      \field{urlmonth}{8}
      \field{urlyear}{2020}
      \field{volume}{29}
      \field{year}{2019}
      \field{urldateera}{ce}
      \field{pages}{89\bibrangedash 129}
      \range{pages}{41}
      \verb{doi}
      \verb 10.1214/18-AAP1403
      \endverb
      \verb{file}
      \verb Nadtochiy and Shkolnikov - 2019 - Particle systems with singular interaction through.pdf:/Users/graemeb/Zotero/storage/E7BEB42T/Nadtochiy and Shkolnikov - 2019 - Particle systems with singular interaction through.pdf:application/pdf
      \endverb
      \verb{urlraw}
      \verb https://projecteuclid.org/euclid.aoap/1544000426
      \endverb
      \verb{url}
      \verb https://projecteuclid.org/euclid.aoap/1544000426
      \endverb
    \endentry
    \entry{nadtochiy_stefan_2022}{misc}{}
      \name{author}{2}{}{%
        {{hash=23bde0d98049dfa33360fa0880afaec5}{%
           family={Nadtochiy},
           familyi={N\bibinitperiod},
           given={Sergey},
           giveni={S\bibinitperiod}}}%
        {{hash=cce28fd978c3edebbc2aee9c6f7e9b0f}{%
           family={Shkolnikov},
           familyi={S\bibinitperiod},
           given={Mykhaylo},
           giveni={M\bibinitperiod}}}%
      }
      \list{publisher}{1}{%
        {arXiv}%
      }
      \strng{namehash}{a0590b36e4b9977c6746f946d087f6c1}
      \strng{fullhash}{a0590b36e4b9977c6746f946d087f6c1}
      \strng{bibnamehash}{a0590b36e4b9977c6746f946d087f6c1}
      \strng{authorbibnamehash}{a0590b36e4b9977c6746f946d087f6c1}
      \strng{authornamehash}{a0590b36e4b9977c6746f946d087f6c1}
      \strng{authorfullhash}{a0590b36e4b9977c6746f946d087f6c1}
      \field{extraname}{3}
      \field{labelalpha}{NS22}
      \field{sortinit}{N}
      \field{sortinithash}{f7242c3ed3dc50029fca1be76c497c7c}
      \field{labelnamesource}{author}
      \field{labeltitlesource}{shorttitle}
      \field{abstract}{We consider the Stefan problem with surface tension, also known as the Stefan-Gibbs-Thomson problem, in an ambient space of arbitrary dimension. Assuming the radial symmetry of the initial data we introduce a novel "probabilistic" notion of solution, which can accommodate the discontinuities in time (of the radius) of the evolving aggregate. Our main result establishes the global existence of a probabilistic solution satisfying the natural upper bound on the sizes of the discontinuities. Moreover, we prove that the upper bound is sharp in dimensions d{>}2, in the sense that none of the discontinuities in the solution can be decreased in magnitude. The detailed analysis of the discontinuities, via appropriate stochastic representations, differentiates this work from the previous literature on weak solutions to the Stefan problem with surface tension.}
      \field{month}{3}
      \field{note}{arXiv:2203.15113 [math-ph]}
      \field{shorttitle}{Stefan problem with surface tension}
      \field{title}{Stefan problem with surface tension: global existence of physical solutions under radial symmetry}
      \field{urlday}{20}
      \field{urlmonth}{6}
      \field{urlyear}{2023}
      \field{year}{2022}
      \field{urldateera}{ce}
      \verb{doi}
      \verb 10.48550/arXiv.2203.15113
      \endverb
      \verb{file}
      \verb arXiv Fulltext PDF:/Users/graemeb/Zotero/storage/E6QGT3WE/Nadtochiy and Shkolnikov - 2022 - Stefan problem with surface tension global existe.pdf:application/pdf;arXiv.org Snapshot:/Users/graemeb/Zotero/storage/B94PF4N5/2203.html:text/html
      \endverb
      \verb{urlraw}
      \verb http://arxiv.org/abs/2203.15113
      \endverb
      \verb{url}
      \verb http://arxiv.org/abs/2203.15113
      \endverb
      \keyw{80A22,60H30,35B44,Mathematical Physics,Mathematics - Analysis of PDEs,Mathematics - Probability}
    \endentry
    \entry{whitt_stochastic-process_2002}{book}{}
      \name{author}{1}{}{%
        {{hash=64244c71a12117d5fe7817ef943cb53a}{%
           family={Whitt},
           familyi={W\bibinitperiod},
           given={Ward},
           giveni={W\bibinitperiod}}}%
      }
      \list{language}{1}{%
        {en}%
      }
      \list{location}{1}{%
        {New York}%
      }
      \list{publisher}{1}{%
        {Springer}%
      }
      \strng{namehash}{64244c71a12117d5fe7817ef943cb53a}
      \strng{fullhash}{64244c71a12117d5fe7817ef943cb53a}
      \strng{bibnamehash}{64244c71a12117d5fe7817ef943cb53a}
      \strng{authorbibnamehash}{64244c71a12117d5fe7817ef943cb53a}
      \strng{authornamehash}{64244c71a12117d5fe7817ef943cb53a}
      \strng{authorfullhash}{64244c71a12117d5fe7817ef943cb53a}
      \field{labelalpha}{Whi02}
      \field{sortinit}{W}
      \field{sortinithash}{ecb89ff85896a47dc313960773ac311d}
      \field{labelnamesource}{author}
      \field{labeltitlesource}{shorttitle}
      \field{isbn}{978-0-387-95358-8}
      \field{series}{Springer series in operations research}
      \field{shorttitle}{Stochastic-process limits}
      \field{title}{Stochastic-process limits: an introduction to stochastic-process limits and their application to queues}
      \field{year}{2002}
      \verb{file}
      \verb Whitt - 2002 - Stochastic-process limits an introduction to stoc.pdf:/Users/graemeb/Zotero/storage/VIIR4MQN/Whitt - 2002 - Stochastic-process limits an introduction to stoc.pdf:application/pdf
      \endverb
      \keyw{Queuing theory,Stochastic processes}
    \endentry
  \enddatalist
\endrefsection
\endinput

